
\documentclass[11pt,a4paper]{article}
\usepackage[a4paper,margin=2.4cm,headheight=15pt]{geometry}
\usepackage{fontspec}
\setmainfont{Segoe UI}
\setsansfont{Segoe UI}
\usepackage{graphicx}
\usepackage{xcolor}
\usepackage{booktabs}
\usepackage{longtable}
\usepackage{array}
\usepackage{float}
\usepackage{lscape}
\usepackage{hyperref}
\graphicspath{{./}{figures/}{../figures/}}
\definecolor{Ink}{HTML}{22272B}
\definecolor{Muted}{HTML}{69737A}
\definecolor{Accent}{HTML}{A64F4F}
\definecolor{Rule}{HTML}{D9D6CE}
\definecolor{Paper}{HTML}{FBFAF7}
\hypersetup{
  colorlinks=true,
  linkcolor=Accent,
  urlcolor=Accent,
  citecolor=Accent,
  pdftitle={Selección, caracterización y requerimientos para un sistema analítico biomédico},
  pdfauthor={Monica Cueto Tapia}
}
\renewcommand{\contentsname}{Tabla de contenidos}
\renewcommand{\listfigurename}{Índice de figuras}
\renewcommand{\listtablename}{Índice de tablas}
\renewcommand{\figurename}{Figura}
\renewcommand{\tablename}{Tabla}
\renewcommand{\refname}{Bibliografía}
\makeatletter
\def\ps@academic{%
  \def\@oddhead{\small Proyecto de Ciencia de Datos\hfil UPEA}%
  \def\@evenhead{\small UPEA\hfil Proyecto de Ciencia de Datos}%
  \def\@oddfoot{\hfil\thepage\hfil}%
  \def\@evenfoot{\hfil\thepage\hfil}%
}
\makeatother
\pagestyle{academic}
\setlength{\parindent}{0pt}
\setlength{\parskip}{0.72em}

\begin{document}

\begin{titlepage}
\thispagestyle{empty}
\pagecolor{Paper}
\begin{center}
\vspace*{1.2cm}
{\Large \textcolor{Accent}{Informe científico reproducible}}\\[1.2cm]
{\Huge \bfseries Selección, caracterización y requerimientos para un sistema analítico biomédico}\\[0.45cm]
{\Large Breast Cancer Wisconsin Diagnostic Dataset}\\[1.6cm]
\rule{0.78\textwidth}{0.6pt}\\[1.2cm]
{\Large \bfseries Monica Cueto Tapia}\\[0.5cm]
{\large Maestría en Inteligencia Artificial, Machine Learning y Data Science}\\[0.25cm]
{\large Universidad Pública de El Alto (UPEA)}\\[1.2cm]
{\large 7 de julio de 2026}\\[2.1cm]
\begin{minipage}{0.78\textwidth}
\small
Documento académico preparado como evidencia de un proyecto reproducible de Ciencia de Datos. El informe conserva los análisis, tablas, gráficos e interpretaciones del reporte HTML original, reorganizados en formato de investigación con numeración, referencias cruzadas, índices y anexos.
\end{minipage}
\vfill
{\small Sin APIs, sin nube, sin descargas de datasets y sin modelos externos de inteligencia artificial.}
\end{center}
\end{titlepage}
\nopagecolor

\pagenumbering{roman}
\tableofcontents
\clearpage
\listoffigures
\clearpage
\listoftables
\clearpage

\pagenumbering{arabic}
\section{Resumen ejecutivo}
Se evaluaron tres datasets disponibles localmente en scikit-learn. El dataset \textbf{Breast Cancer Wisconsin Diagnostic} obtuvo el mayor puntaje ponderado (97.7355/100) al combinar tamaño muestral, riqueza dimensional, completitud, preparación analítica, interpretabilidad e impacto. La selección es apropiada para una práctica de maestría porque permite discutir calidad de datos, separación estadística entre clases, redundancia de variables y requerimientos de un flujo analítico reproducible sin depender de servicios externos.

El conjunto seleccionado contiene 569 observaciones, 30 variables predictoras y 0\% de valores faltantes. La evidencia univariada más fuerte aparece en \textit{worst concave points} (|d|=2.6926), mientras que el análisis de correlación evidencia redundancia fuerte entre medidas geométricas. El tablero ejecutivo de la Figura~\ref{fig:dashboard} sintetiza estos hallazgos.


\begin{figure}[H]
\centering
\includegraphics[width=\linewidth,height=0.62\textheight,keepaspectratio]{figures/12_executive_dashboard.png}
\caption{Dashboard ejecutivo del dataset seleccionado}
\label{fig:dashboard}
\vspace{0.2em}
\begin{minipage}{0.92\linewidth}
\footnotesize \textbf{Interpretación.} El panel resume que el conjunto es completo, moderadamente desbalanceado y con variables de fuerte separación diagnóstica.
\end{minipage}
\end{figure}


\section{Fase 1. Búsqueda, comparación y selección del dataset}
Para cumplir la restricción de no usar internet, APIs ni servicios en la nube, la búsqueda se limitó al catálogo local de datasets incluidos en scikit-learn. Esta decisión reduce el consumo computacional, elimina riesgos de disponibilidad externa y mantiene la reproducibilidad. Los candidatos fueron Breast Cancer Wisconsin Diagnostic, Wine Recognition y Diabetes Progression.

Como se muestra en la Figura~\ref{fig:project-flow}, el flujo separa selección, análisis y comunicación. La comparación cuantitativa se resume en las Tablas~\ref{tab:comparison-summary} y~\ref{tab:comparison-scores}; las Figuras~\ref{fig:comparison-matrix} y~\ref{fig:dataset-radar} complementan la decisión visualmente.


\begin{figure}[H]
\centering
\includegraphics[width=\linewidth,height=0.62\textheight,keepaspectratio]{figures/01_project_flow.png}
\caption{Flujo metodológico del proyecto}
\label{fig:project-flow}
\vspace{0.2em}
\begin{minipage}{0.92\linewidth}
\footnotesize \textbf{Interpretación.} El flujo separa selección, análisis y comunicación para que cada conclusión se origine en evidencia verificable.
\end{minipage}
\end{figure}



\scriptsize
\setlength{\tabcolsep}{1.5pt}
\renewcommand{\arraystretch}{1.16}
\begin{longtable}{>{\raggedright\arraybackslash}p{0.167\linewidth} >{\raggedright\arraybackslash}p{0.176\linewidth} >{\raggedright\arraybackslash}p{0.130\linewidth} >{\raggedright\arraybackslash}p{0.093\linewidth} >{\raggedright\arraybackslash}p{0.102\linewidth} >{\raggedright\arraybackslash}p{0.111\linewidth} >{\raggedright\arraybackslash}p{0.121\linewidth}}
\caption{Comparación técnica resumida de datasets candidatos}\label{tab:comparison-summary}\\
\toprule
Dataset & Dominio & Tarea & Obs. & Variables & Faltantes (\%) & Puntaje total \\
\midrule
\endfirsthead
\caption[]{Comparación técnica resumida de datasets candidatos (continuación)}\\
\toprule
Dataset & Dominio & Tarea & Obs. & Variables & Faltantes (\%) & Puntaje total \\
\midrule
\endhead
\midrule
\multicolumn{7}{r}{Continúa en la siguiente página}\\
\midrule
\endfoot
\bottomrule
\endlastfoot
Breast Cancer Wisconsin Diagnostic & Salud / diagnóstico oncológico & Clasificación binaria & 569.0 & 30.00 & 0 & 97.74 \\
Diabetes Progression & Salud / progresión metabólica & Regresión & 442.0 & 10.00 & 0 & 83.53 \\
Wine Recognition & Química analítica / enología & Clasificación multiclase & 178.0 & 13.00 & 0 & 76.18 \\
\end{longtable}
\normalsize


\begin{landscape}

\scriptsize
\setlength{\tabcolsep}{1.5pt}
\renewcommand{\arraystretch}{1.16}
\begin{longtable}{>{\raggedright\arraybackslash}p{0.154\linewidth} >{\raggedright\arraybackslash}p{0.094\linewidth} >{\raggedright\arraybackslash}p{0.111\linewidth} >{\raggedright\arraybackslash}p{0.103\linewidth} >{\raggedright\arraybackslash}p{0.128\linewidth} >{\raggedright\arraybackslash}p{0.111\linewidth} >{\raggedright\arraybackslash}p{0.077\linewidth} >{\raggedright\arraybackslash}p{0.103\linewidth}}
\caption{Puntajes normalizados usados para la selección del dataset}\label{tab:comparison-scores}\\
\toprule
Dataset & Tamaño & Dimensión & Completitud & Prep. analítica & Interpretabilidad & Impacto & Puntaje total \\
\midrule
\endfirsthead
\caption[]{Puntajes normalizados usados para la selección del dataset (continuación)}\\
\toprule
Dataset & Tamaño & Dimensión & Completitud & Prep. analítica & Interpretabilidad & Impacto & Puntaje total \\
\midrule
\endhead
\midrule
\multicolumn{8}{r}{Continúa en la siguiente página}\\
\midrule
\endfoot
\bottomrule
\endlastfoot
Breast Cancer Wisconsin Diagnostic & 100.0 & 100.0 & 100.0 & 94.90 & 90.00 & 100.0 & 97.74 \\
Diabetes Progression & 88.14 & 57.73 & 100.0 & 100.0 & 73.00 & 90.00 & 83.53 \\
Wine Recognition & 55.93 & 65.83 & 100.0 & 96.18 & 86.00 & 67.00 & 76.18 \\
\end{longtable}
\normalsize

\end{landscape}



\begin{figure}[H]
\centering
\includegraphics[width=\linewidth,height=0.62\textheight,keepaspectratio]{figures/02_dataset_comparison_matrix.png}
\caption{Matriz comparativa de datasets candidatos}
\label{fig:comparison-matrix}
\vspace{0.2em}
\begin{minipage}{0.92\linewidth}
\footnotesize \textbf{Interpretación.} El dataset de cáncer de mama domina por tamaño, profundidad de variables, impacto e interpretabilidad; la completitud es equivalente en los tres.
\end{minipage}
\end{figure}



\begin{figure}[H]
\centering
\includegraphics[width=\linewidth,height=0.62\textheight,keepaspectratio]{figures/03_dataset_radar.png}
\caption{Perfil multicriterio de selección}
\label{fig:dataset-radar}
\vspace{0.2em}
\begin{minipage}{0.92\linewidth}
\footnotesize \textbf{Interpretación.} La decisión no depende de un único criterio: el dataset seleccionado mantiene desempeño alto en seis dimensiones.
\end{minipage}
\end{figure}


La ponderación priorizó evidencia verificable: 25\% tamaño muestral, 20\% dimensionalidad, 15\% completitud, 15\% preparación analítica, 15\% interpretabilidad y 10\% impacto. El dataset seleccionado domina por volumen, número de variables, completitud total y relevancia biomédica. Wine Recognition es técnicamente limpio, pero tiene menor escala e impacto para una discusión socio-técnica. Diabetes Progression posee impacto alto, aunque su objetivo continuo y variables estandarizadas reducen interpretabilidad variable por variable.

\section{Fase 2. Contexto, problema, objetivo y alcance}
\subsection{Contexto}
El análisis biomédico basado en imágenes requiere transformar mediciones morfológicas en evidencia interpretable. Este dataset resume características geométricas y de textura de núcleos celulares obtenidas a partir de aspiración con aguja fina de masas mamarias.

\subsection{Problema}
Antes de plantear un sistema predictivo, es necesario demostrar que los datos son completos, comprensibles, estadísticamente informativos y gobernables. Sin esa base, cualquier modelo posterior podría amplificar redundancias, sesgos de clase o interpretaciones clínicas débiles.

\subsection{Objetivo}
Seleccionar y caracterizar un dataset biomédico reproducible, identificar sus variables críticas, evaluar calidad estadística y definir requerimientos de un sistema analítico que pueda sostener futuras fases de modelado con rigor académico.

\subsection{Alcance}
El alcance cubre selección de dataset, análisis exploratorio, diccionario de datos, visualizaciones editoriales y requerimientos. No incluye entrenamiento de modelos clínicos ni recomendaciones médicas individuales, porque las fases solicitadas no lo exigen y porque se prioriza eficiencia.

\subsection{Beneficiarios}
Los beneficiarios son estudiantes e investigadores de ciencia de datos, docentes evaluadores, equipos de analítica biomédica y responsables de gobernanza de datos que requieran una base reproducible para decisiones metodológicas.

\subsection{Resultados esperados}
Se espera una selección defendible, un inventario de variables accionable, evidencia de calidad de datos, detección de señales estadísticas relevantes, visualizaciones publicables y especificación clara de requerimientos para continuidad del proyecto.


\small
\setlength{\tabcolsep}{1.5pt}
\renewcommand{\arraystretch}{1.16}
\begin{longtable}{>{\raggedright\arraybackslash}p{0.191\linewidth} >{\raggedright\arraybackslash}p{0.201\linewidth} >{\raggedright\arraybackslash}p{0.507\linewidth}}
\caption{Resumen de calidad del dataset seleccionado}\label{tab:quality}\\
\toprule
Indicador & Resultado & Interpretación técnica \\
\midrule
\endfirsthead
\caption[]{Resumen de calidad del dataset seleccionado (continuación)}\\
\toprule
Indicador & Resultado & Interpretación técnica \\
\midrule
\endhead
\midrule
\multicolumn{3}{r}{Continúa en la siguiente página}\\
\midrule
\endfoot
\bottomrule
\endlastfoot
Observaciones & 569.0 & Tamaño suficiente para exploración estadística y especificación de requerimientos. \\
Variables predictoras & 30.00 & Cobertura morfológica amplia: medias, errores estándar y peores valores. \\
Variable objetivo & diagnosis & Etiqueta clínica binaria para comparar patrones entre masas benignas y malignas. \\
Valores faltantes & 0 (0.00\%) & No se requiere imputación; se preserva trazabilidad. \\
Filas duplicadas & 0 & No se detectan registros duplicados exactos en el conjunto final. \\
Distribución objetivo & benign: 357; malignant: 212 & Existe desbalance moderado, no extremo. \\
Proporción minoritaria & 37.3\% & La clase maligna conserva representación analítica suficiente. \\
Variables con atípicos IQR & 29 de 30 & Los atípicos deben revisarse como extremos biológicos plausibles, no eliminarse por defecto. \\
Mayor correlación absoluta & mean radius vs mean perimeter: 0.998 & Sugiere redundancia morfológica que debe controlarse en modelado posterior. \\
Variables con |d| >= 1 & 16.00 & La separación univariada es fuerte en varias variables, en especial medidas de concavidad y tamaño. \\
\end{longtable}
\normalsize


La Figura~\ref{fig:missingness} confirma completitud total; la Figura~\ref{fig:target-distribution} muestra un desbalance moderado, no extremo. La señal univariada se observa en la Figura~\ref{fig:effect-size} y se complementa con la Tabla~\ref{tab:effect-sizes}.


\begin{figure}[H]
\centering
\includegraphics[width=\linewidth,height=0.62\textheight,keepaspectratio]{figures/04_missingness_map.png}
\caption{Diagrama de completitud por familia de variable}
\label{fig:missingness}
\vspace{0.2em}
\begin{minipage}{0.92\linewidth}
\footnotesize \textbf{Interpretación.} La completitud uniforme del 100\% elimina la necesidad de imputación en esta fase y reduce incertidumbre por datos faltantes.
\end{minipage}
\end{figure}



\begin{figure}[H]
\centering
\includegraphics[width=\linewidth,height=0.62\textheight,keepaspectratio]{figures/05_target_distribution.png}
\caption{Distribución de la variable objetivo}
\label{fig:target-distribution}
\vspace{0.2em}
\begin{minipage}{0.92\linewidth}
\footnotesize \textbf{Interpretación.} La clase benigna es mayoritaria, pero la clase maligna representa más de un tercio del conjunto; el desbalance es manejable para análisis exploratorio.
\end{minipage}
\end{figure}



\begin{figure}[H]
\centering
\includegraphics[width=\linewidth,height=0.62\textheight,keepaspectratio]{figures/06_effect_size_lollipop.png}
\caption{Variables con mayor separación entre diagnósticos}
\label{fig:effect-size}
\vspace{0.2em}
\begin{minipage}{0.92\linewidth}
\footnotesize \textbf{Interpretación.} Las medidas de concavidad, puntos cóncavos y perímetro/área concentran la mayor diferencia estandarizada entre masas malignas y benignas.
\end{minipage}
\end{figure}


\begin{landscape}

\scriptsize
\setlength{\tabcolsep}{1.5pt}
\renewcommand{\arraystretch}{1.16}
\begin{longtable}{>{\raggedright\arraybackslash}p{0.200\linewidth} >{\raggedright\arraybackslash}p{0.130\linewidth} >{\raggedright\arraybackslash}p{0.130\linewidth} >{\raggedright\arraybackslash}p{0.160\linewidth} >{\raggedright\arraybackslash}p{0.090\linewidth} >{\raggedright\arraybackslash}p{0.170\linewidth}}
\caption{Tamaños de efecto por variable entre diagnósticos}\label{tab:effect-sizes}\\
\toprule
Variable & Media maligno & Media benigno & Diferencia estandarizada (d) & |d| & Dirección \\
\midrule
\endfirsthead
\caption[]{Tamaños de efecto por variable entre diagnósticos (continuación)}\\
\toprule
Variable & Media maligno & Media benigno & Diferencia estandarizada (d) & |d| & Dirección \\
\midrule
\endhead
\midrule
\multicolumn{6}{r}{Continúa en la siguiente página}\\
\midrule
\endfoot
\bottomrule
\endlastfoot
worst concave points & 0.1822 & 0.0744 & 2.6926 & 2.6926 & Mayor en maligno \\
worst perimeter & 141.4 & 87.01 & 2.5982 & 2.5982 & Mayor en maligno \\
mean concave points & 0.088 & 0.0257 & 2.5452 & 2.5452 & Mayor en maligno \\
worst radius & 21.13 & 13.38 & 2.5439 & 2.5439 & Mayor en maligno \\
mean perimeter & 115.4 & 78.08 & 2.2895 & 2.2895 & Mayor en maligno \\
worst area & 1,422.3 & 558.9 & 2.2302 & 2.2302 & Mayor en maligno \\
mean radius & 17.46 & 12.15 & 2.2055 & 2.2055 & Mayor en maligno \\
mean area & 978.4 & 462.8 & 2.0757 & 2.0757 & Mayor en maligno \\
mean concavity & 0.1608 & 0.0461 & 2.0033 & 2.0033 & Mayor en maligno \\
worst concavity & 0.4506 & 0.1662 & 1.8119 & 1.8119 & Mayor en maligno \\
mean compactness & 0.1452 & 0.0801 & 1.5346 & 1.5346 & Mayor en maligno \\
worst compactness & 0.3748 & 0.1827 & 1.5126 & 1.5126 & Mayor en maligno \\
radius error & 0.6091 & 0.2841 & 1.4217 & 1.4217 & Mayor en maligno \\
perimeter error & 4.3239 & 2.0003 & 1.3816 & 1.3816 & Mayor en maligno \\
area error & 72.67 & 21.14 & 1.3534 & 1.3534 & Mayor en maligno \\
worst texture & 29.32 & 23.52 & 1.0605 & 1.0605 & Mayor en maligno \\
worst smoothness & 0.1448 & 0.125 & 0.9596 & 0.9596 & Mayor en maligno \\
worst symmetry & 0.3235 & 0.2702 & 0.9453 & 0.9453 & Mayor en maligno \\
mean texture & 21.60 & 17.91 & 0.9423 & 0.9423 & Mayor en maligno \\
concave points error & 0.0151 & 0.0099 & 0.9228 & 0.9228 & Mayor en maligno \\
mean smoothness & 0.1029 & 0.0925 & 0.793 & 0.793 & Mayor en maligno \\
mean symmetry & 0.1929 & 0.1742 & 0.723 & 0.723 & Mayor en maligno \\
worst fractal dimension & 0.0915 & 0.0794 & 0.7068 & 0.7068 & Mayor en maligno \\
compactness error & 0.0323 & 0.0214 & 0.6327 & 0.6327 & Mayor en maligno \\
concavity error & 0.0418 & 0.026 & 0.5416 & 0.5416 & Mayor en maligno \\
fractal dimension error & 0.0041 & 0.0036 & 0.1615 & 0.1615 & Mayor en maligno \\
smoothness error & 0.0068 & 0.0072 & -0.1387 & 0.1387 & Mayor en benigno \\
mean fractal dimension & 0.0627 & 0.0629 & -0.0265 & 0.0265 & Mayor en benigno \\
texture error & 1.2109 & 1.2204 & -0.0171 & 0.0171 & Mayor en benigno \\
symmetry error & 0.0205 & 0.0206 & -0.0135 & 0.0135 & Mayor en benigno \\
\end{longtable}
\normalsize

\end{landscape}


La Figura~\ref{fig:correlation} evidencia que la señal predictiva convive con multicolinealidad potencial. Las Figuras~\ref{fig:distributions} y~\ref{fig:boxplots} muestran los desplazamientos de distribución y mediana entre clases para las variables prioritarias.


\begin{figure}[H]
\centering
\includegraphics[width=\linewidth,height=0.62\textheight,keepaspectratio]{figures/07_correlation_heatmap.png}
\caption{Heatmap de correlación entre variables de mayor señal}
\label{fig:correlation}
\vspace{0.2em}
\begin{minipage}{0.92\linewidth}
\footnotesize \textbf{Interpretación.} La señal predictiva convive con redundancia fuerte entre medidas geométricas; esto sugiere cautela con multicolinealidad en fases de modelado.
\end{minipage}
\end{figure}



\begin{figure}[H]
\centering
\includegraphics[width=\linewidth,height=0.62\textheight,keepaspectratio]{figures/08_distribution_panel.png}
\caption{Distribuciones comparadas de variables prioritarias}
\label{fig:distributions}
\vspace{0.2em}
\begin{minipage}{0.92\linewidth}
\footnotesize \textbf{Interpretación.} Las distribuciones muestran desplazamientos sistemáticos entre clases, especialmente en variables de borde y concavidad.
\end{minipage}
\end{figure}



\begin{figure}[H]
\centering
\includegraphics[width=\linewidth,height=0.62\textheight,keepaspectratio]{figures/09_boxplot_panel.png}
\caption{Boxplots robustos por diagnóstico}
\label{fig:boxplots}
\vspace{0.2em}
\begin{minipage}{0.92\linewidth}
\footnotesize \textbf{Interpretación.} La escala robusta permite comparar variables con unidades diferentes y evidencia medianas más altas en malignidad para variables clave.
\end{minipage}
\end{figure}


\section{Fase 3. Requerimientos del sistema analítico}
Los requerimientos se formulan para un sistema de análisis reproducible, no para un producto clínico operativo. La prioridad se asigna según riesgo metodológico, trazabilidad y valor para decisiones académicas posteriores. La arquitectura propuesta se observa en la Figura~\ref{fig:architecture} y el cronograma ejecutivo en la Figura~\ref{fig:timeline}.


\begin{figure}[H]
\centering
\includegraphics[width=\linewidth,height=0.62\textheight,keepaspectratio]{figures/10_architecture_pipeline.png}
\caption{Arquitectura reproducible del proyecto}
\label{fig:architecture}
\vspace{0.2em}
\begin{minipage}{0.92\linewidth}
\footnotesize \textbf{Interpretación.} La arquitectura mantiene una sola ruta de datos y separa artefactos finales por función, evitando duplicación y procesamiento innecesario.
\end{minipage}
\end{figure}



\begin{figure}[H]
\centering
\includegraphics[width=\linewidth,height=0.62\textheight,keepaspectratio]{figures/11_project_timeline.png}
\caption{Cronograma ejecutivo de la práctica}
\label{fig:timeline}
\vspace{0.2em}
\begin{minipage}{0.92\linewidth}
\footnotesize \textbf{Interpretación.} El cronograma prioriza selección y comprensión de datos antes de cualquier modelado, consistente con una práctica de formulación analítica.
\end{minipage}
\end{figure}


\begin{landscape}

\scriptsize
\setlength{\tabcolsep}{1.5pt}
\renewcommand{\arraystretch}{1.16}
\begin{longtable}{>{\raggedright\arraybackslash}p{0.073\linewidth} >{\raggedright\arraybackslash}p{0.266\linewidth} >{\raggedright\arraybackslash}p{0.119\linewidth} >{\raggedright\arraybackslash}p{0.083\linewidth} >{\raggedright\arraybackslash}p{0.193\linewidth} >{\raggedright\arraybackslash}p{0.147\linewidth}}
\caption{Requerimientos funcionales}\label{tab:functional}\\
\toprule
ID & Requerimiento funcional & Categoría & Prioridad & Justificación & Impacto \\
\midrule
\endfirsthead
\caption[]{Requerimientos funcionales (continuación)}\\
\toprule
ID & Requerimiento funcional & Categoría & Prioridad & Justificación & Impacto \\
\midrule
\endhead
\midrule
\multicolumn{6}{r}{Continúa en la siguiente página}\\
\midrule
\endfoot
\bottomrule
\endlastfoot
RF-01 & Cargar el dataset seleccionado desde una fuente local versionada. & Ingesta & Alta & Evita dependencias externas y garantiza reproducción en entornos académicos. & Permite reconstruir el análisis sin internet ni APIs. \\
RF-02 & Validar estructura, número de variables, tipos y variable objetivo. & Calidad de datos & Alta & Previene análisis sobre columnas corruptas o incompatibles. & Reduce riesgo de conclusiones no trazables. \\
RF-03 & Calcular métricas de completitud, duplicados y balance del objetivo. & Calidad de datos & Alta & La confiabilidad del diagnóstico depende de ausencia de sesgos de captura. & Define si se requiere imputación, depuración o estratificación. \\
RF-04 & Generar comparación cuantitativa entre tres datasets candidatos. & Selección analítica & Alta & La selección debe sustentarse en criterios medibles y no en preferencia subjetiva. & Justifica el dataset ante evaluación académica. \\
RF-05 & Construir diccionario exhaustivo de variables con rango, unidad e importancia. & Documentación & Alta & La interpretación clínica necesita semántica clara por variable. & Mejora auditabilidad y transferencia del proyecto. \\
RF-06 & Calcular estadísticos descriptivos robustos por variable. & EDA & Alta & Percentiles, asimetría y atípicos revelan la estructura real de los datos. & Soporta decisiones de transformación y modelado futuro. \\
RF-07 & Cuantificar separación entre clases mediante tamaño de efecto. & Análisis estadístico & Alta & El tamaño de efecto aporta evidencia interpretable sin entrenar modelos innecesarios. & Prioriza variables relevantes para diagnóstico. \\
RF-08 & Identificar correlaciones altas entre variables predictoras. & Análisis estadístico & Media & La redundancia puede afectar interpretabilidad y estabilidad de modelos futuros. & Informa reducción dimensional o regularización posterior. \\
RF-09 & Exportar tablas finales en formato reutilizable. & Entregables & Media & Las tablas deben poder revisarse, auditarse o integrarse en anexos. & Facilita evaluación y continuidad del trabajo. \\
RF-10 & Compilar un reporte HTML editorial con figuras, tablas e interpretación. & Reporte & Alta & El entregable final debe comunicar evidencia, no solo producir archivos. & Aumenta claridad, profesionalismo y presentabilidad. \\
\end{longtable}
\normalsize

\end{landscape}


\begin{landscape}

\scriptsize
\setlength{\tabcolsep}{1.5pt}
\renewcommand{\arraystretch}{1.16}
\begin{longtable}{>{\raggedright\arraybackslash}p{0.084\linewidth} >{\raggedright\arraybackslash}p{0.206\linewidth} >{\raggedright\arraybackslash}p{0.094\linewidth} >{\raggedright\arraybackslash}p{0.281\linewidth} >{\raggedright\arraybackslash}p{0.215\linewidth}}
\caption{Requerimientos no funcionales}\label{tab:nonfunctional}\\
\toprule
ID & Requerimiento no funcional & Prioridad & Justificación & Impacto \\
\midrule
\endfirsthead
\caption[]{Requerimientos no funcionales (continuación)}\\
\toprule
ID & Requerimiento no funcional & Prioridad & Justificación & Impacto \\
\midrule
\endhead
\midrule
\multicolumn{5}{r}{Continúa en la siguiente página}\\
\midrule
\endfoot
\bottomrule
\endlastfoot
RNF-01 & Reproducibilidad & Alta & El flujo debe ejecutarse con un único script local y resultados deterministas. & Permite validar el trabajo en cualquier revisión académica. \\
RNF-02 & Eficiencia computacional & Alta & Solo se cargan datasets livianos y se evita entrenamiento innecesario. & Minimiza tiempo, memoria y consumo de recursos. \\
RNF-03 & Independencia de red & Alta & No se usan descargas, APIs, nube ni modelos externos. & El proyecto cumple restricciones y protege continuidad operativa. \\
RNF-04 & Trazabilidad & Alta & Cada figura y tabla declara fuente e interpretación. & Facilita auditoría de evidencia y defensa metodológica. \\
RNF-05 & Interpretabilidad & Alta & Las métricas se basan en estadística descriptiva y tamaño de efecto. & Evita cajas negras en fases tempranas del proyecto. \\
RNF-06 & Mantenibilidad & Media & El código separa carga, métricas, visualización y reporte. & Permite extender el análisis sin rehacer la base. \\
RNF-07 & Portabilidad & Media & Usa únicamente Python, Pandas, NumPy, Matplotlib y scikit-learn. & Reduce fricción de instalación y dependencia de librerías pesadas. \\
RNF-08 & Calidad visual & Alta & La salida debe usar paleta sobria, espacio en blanco y jerarquía editorial. & Mejora comunicación de hallazgos a audiencias técnicas y ejecutivas. \\
RNF-09 & Robustez & Media & Las validaciones deben fallar de forma explícita ante cambios de estructura. & Evita reportes silenciosamente inconsistentes. \\
RNF-10 & Gobernanza ética & Alta & El análisis evita inferencias clínicas individualizadas y trata el dataset como material académico. & Reduce riesgo de uso indebido en decisiones médicas reales. \\
\end{longtable}
\normalsize

\end{landscape}


\section{Conclusiones}
\begin{itemize}
\item La selección del dataset es cuantitativamente defendible: Breast Cancer Wisconsin Diagnostic alcanza 97.7355/100 y supera a los candidatos por balance entre escala, dimensionalidad, interpretabilidad e impacto.
\item La completitud del 100\% evita imputación en esta fase; por tanto, las diferencias observadas entre clases no están condicionadas por estrategias de reemplazo de datos faltantes.
\item La clase maligna conserva representación analítica suficiente para exploración, aunque un modelado posterior debería controlar el desbalance con particiones estratificadas.
\item La presencia de múltiples variables con |d| $\geq$ 1 indica separación estadística fuerte; no obstante, la correlación alta entre medidas geométricas confirma información redundante.
\item La fase de requerimientos debe priorizar reproducibilidad, trazabilidad, independencia de red y comunicación visual, porque esos atributos hacen defendible el proyecto en un contexto de maestría.
\end{itemize}

\section{Bibliografía}
\begin{thebibliography}{9}
\bibitem{pedregosa2011} Pedregosa, F. et al. (2011). \textit{Scikit-learn: Machine Learning in Python}. Journal of Machine Learning Research, 12, 2825--2830.
\bibitem{street1993} Street, W. N., Wolberg, W. H., \& Mangasarian, O. L. (1993). \textit{Nuclear feature extraction for breast tumor diagnosis}. IS\&T/SPIE International Symposium on Electronic Imaging.
\bibitem{wolberg1995} Wolberg, W. H., Street, W. N., \& Mangasarian, O. L. (1995). \textit{Breast cancer diagnosis and prognosis via linear programming}. Operations Research, 43(4), 570--577.
\end{thebibliography}

\clearpage
\appendix
\section{Anexo A. Diccionario exhaustivo de variables}
La Tabla~\ref{tab:data-dictionary} reproduce el diccionario completo de variables del proyecto, con tipo, descripción, rango observado, unidad, importancia, posibles valores atípicos y observaciones técnicas.

\begin{landscape}

\tiny
\setlength{\tabcolsep}{1.5pt}
\renewcommand{\arraystretch}{1.16}
\begin{longtable}{>{\raggedright\arraybackslash}p{0.101\linewidth} >{\raggedright\arraybackslash}p{0.084\linewidth} >{\raggedright\arraybackslash}p{0.168\linewidth} >{\raggedright\arraybackslash}p{0.084\linewidth} >{\raggedright\arraybackslash}p{0.101\linewidth} >{\raggedright\arraybackslash}p{0.075\linewidth} >{\raggedright\arraybackslash}p{0.134\linewidth} >{\raggedright\arraybackslash}p{0.134\linewidth}}
\caption{Diccionario exhaustivo de variables}\label{tab:data-dictionary}\\
\toprule
Nombre & Tipo & Descripción & Rango / valores & Unidad & Importancia & Posibles valores atípicos & Obs. \\
\midrule
\endfirsthead
\caption[]{Diccionario exhaustivo de variables (continuación)}\\
\toprule
Nombre & Tipo & Descripción & Rango / valores & Unidad & Importancia & Posibles valores atípicos & Obs. \\
\midrule
\endhead
\midrule
\multicolumn{8}{r}{Continúa en la siguiente página}\\
\midrule
\endfoot
\bottomrule
\endlastfoot
diagnosis & categórica nominal & Variable objetivo: diagnóstico histopatológico codificado como maligno o benigno. & malignant, benign & categoría clínica & Crítica & No aplica; se valida consistencia de etiquetas. & Es la variable de referencia para segmentar el análisis estadístico. \\
mean radius & numérica continua & Media de distancia media desde el centro hasta puntos del perímetro celular. & 6.981 a 28.11 & unidades de imagen & Muy alta & 14 casos (2.5\%) fuera de [5.58, 21.9] por regla IQR. & mayor en maligno; |d|=2.21; asimetría moderada. \\
mean texture & numérica continua & Media de variabilidad de la intensidad de gris en la imagen. & 9.71 a 39.28 & desviación estándar de niveles de gris & Media & 7 casos (1.2\%) fuera de [7.73, 30.2] por regla IQR. & mayor en maligno; |d|=0.94; asimetría moderada. \\
mean perimeter & numérica continua & Media de longitud del contorno de la masa celular. & 43.79 a 188.5 & unidades de imagen & Muy alta & 13 casos (2.3\%) fuera de [31.8, 147] por regla IQR. & mayor en maligno; |d|=2.29; asimetría moderada. \\
mean area & numérica continua & Media de superficie proyectada de la masa celular. & 143.5 a 2501 & unidades cuadradas de imagen & Muy alta & 25 casos (4.4\%) fuera de [-123, 1.33e+03] por regla IQR. & mayor en maligno; |d|=2.08; asimetría marcada. \\
mean smoothness & numérica continua & Media de variación local de longitudes de radio. & 0.05263 a 0.1634 & índice adimensional & Media & 6 casos (1.1\%) fuera de [0.058, 0.134] por regla IQR. & mayor en maligno; |d|=0.79; distribución relativamente estable. \\
mean compactness & numérica continua & Media de relación de compacidad calculada como perímetro al cuadrado sobre área menos uno. & 0.01938 a 0.3454 & índice adimensional & Alta & 16 casos (2.8\%) fuera de [-0.0333, 0.229] por regla IQR. & mayor en maligno; |d|=1.53; asimetría marcada. \\
mean concavity & numérica continua & Media de severidad de porciones cóncavas del contorno. & 0 a 0.4268 & índice adimensional & Muy alta & 18 casos (3.2\%) fuera de [-0.122, 0.282] por regla IQR. & mayor en maligno; |d|=2.00; asimetría marcada. \\
mean concave points & numérica continua & Media de cantidad relativa de puntos cóncavos del contorno. & 0 a 0.2012 & índice adimensional & Muy alta & 10 casos (1.8\%) fuera de [-0.0602, 0.155] por regla IQR. & mayor en maligno; |d|=2.55; asimetría marcada. \\
mean symmetry & numérica continua & Media de grado de simetría de la estructura celular. & 0.106 a 0.304 & índice adimensional & Media & 15 casos (2.6\%) fuera de [0.111, 0.246] por regla IQR. & mayor en maligno; |d|=0.72; asimetría moderada. \\
mean fractal dimension & numérica continua & Media de aproximación de complejidad del contorno celular. & 0.04996 a 0.09744 & índice adimensional & Exploratoria & 15 casos (2.6\%) fuera de [0.0451, 0.0787] por regla IQR. & mayor en benigno; |d|=0.03; asimetría marcada. \\
radius error & numérica continua & Error estándar de distancia media desde el centro hasta puntos del perímetro celular. & 0.1115 a 2.873 & unidades de imagen & Alta & 38 casos (6.7\%) fuera de [-0.137, 0.849] por regla IQR. & mayor en maligno; |d|=1.42; asimetría marcada. \\
texture error & numérica continua & Error estándar de variabilidad de la intensidad de gris en la imagen. & 0.3602 a 4.885 & desviación estándar de niveles de gris & Exploratoria & 20 casos (3.5\%) fuera de [-0.126, 2.43] por regla IQR. & mayor en benigno; |d|=0.02; asimetría marcada. \\
perimeter error & numérica continua & Error estándar de longitud del contorno de la masa celular. & 0.757 a 21.98 & unidades de imagen & Alta & 38 casos (6.7\%) fuera de [-1.02, 5.98] por regla IQR. & mayor en maligno; |d|=1.38; asimetría marcada. \\
area error & numérica continua & Error estándar de superficie proyectada de la masa celular. & 6.802 a 542.2 & unidades cuadradas de imagen & Alta & 65 casos (11.4\%) fuera de [-23.2, 86.2] por regla IQR. & mayor en maligno; |d|=1.35; asimetría marcada. \\
smoothness error & numérica continua & Error estándar de variación local de longitudes de radio. & 0.001713 a 0.03113 & índice adimensional & Exploratoria & 30 casos (5.3\%) fuera de [0.000703, 0.0126] por regla IQR. & mayor en benigno; |d|=0.14; asimetría marcada. \\
compactness error & numérica continua & Error estándar de relación de compacidad calculada como perímetro al cuadrado sobre área menos uno. & 0.002252 a 0.1354 & índice adimensional & Exploratoria & 28 casos (4.9\%) fuera de [-0.016, 0.0615] por regla IQR. & mayor en maligno; |d|=0.63; asimetría marcada. \\
concavity error & numérica continua & Error estándar de severidad de porciones cóncavas del contorno. & 0 a 0.396 & índice adimensional & Exploratoria & 22 casos (3.9\%) fuera de [-0.0253, 0.0825] por regla IQR. & mayor en maligno; |d|=0.54; asimetría marcada. \\
concave points error & numérica continua & Error estándar de cantidad relativa de puntos cóncavos del contorno. & 0 a 0.05279 & índice adimensional & Media & 19 casos (3.3\%) fuera de [-0.00297, 0.0253] por regla IQR. & mayor en maligno; |d|=0.92; asimetría marcada. \\
symmetry error & numérica continua & Error estándar de grado de simetría de la estructura celular. & 0.007882 a 0.07895 & índice adimensional & Exploratoria & 27 casos (4.7\%) fuera de [0.00268, 0.036] por regla IQR. & mayor en benigno; |d|=0.01; asimetría marcada. \\
fractal dimension error & numérica continua & Error estándar de aproximación de complejidad del contorno celular. & 0.0008948 a 0.02984 & índice adimensional & Exploratoria & 28 casos (4.9\%) fuera de [-0.00122, 0.00802] por regla IQR. & mayor en maligno; |d|=0.16; asimetría marcada. \\
worst radius & numérica continua & Peor valor de distancia media desde el centro hasta puntos del perímetro celular. & 7.93 a 36.04 & unidades de imagen & Muy alta & 17 casos (3.0\%) fuera de [4.34, 27.5] por regla IQR. & mayor en maligno; |d|=2.54; asimetría marcada. \\
worst texture & numérica continua & Peor valor de variabilidad de la intensidad de gris en la imagen. & 12.02 a 49.54 & desviación estándar de niveles de gris & Media & 5 casos (0.9\%) fuera de [8.12, 42.7] por regla IQR. & mayor en maligno; |d|=1.06; distribución relativamente estable. \\
worst perimeter & numérica continua & Peor valor de longitud del contorno de la masa celular. & 50.41 a 251.2 & unidades de imagen & Muy alta & 15 casos (2.6\%) fuera de [22.2, 187] por regla IQR. & mayor en maligno; |d|=2.60; asimetría marcada. \\
worst area & numérica continua & Peor valor de superficie proyectada de la masa celular. & 185.2 a 4254 & unidades cuadradas de imagen & Muy alta & 35 casos (6.2\%) fuera de [-338, 1.94e+03] por regla IQR. & mayor en maligno; |d|=2.23; asimetría marcada. \\
worst smoothness & numérica continua & Peor valor de variación local de longitudes de radio. & 0.07117 a 0.2226 & índice adimensional & Media & 7 casos (1.2\%) fuera de [0.0725, 0.19] por regla IQR. & mayor en maligno; |d|=0.96; distribución relativamente estable. \\
worst compactness & numérica continua & Peor valor de relación de compacidad calculada como perímetro al cuadrado sobre área menos uno. & 0.02729 a 1.058 & índice adimensional & Alta & 16 casos (2.8\%) fuera de [-0.141, 0.627] por regla IQR. & mayor en maligno; |d|=1.51; asimetría marcada. \\
worst concavity & numérica continua & Peor valor de severidad de porciones cóncavas del contorno. & 0 a 1.252 & índice adimensional & Alta & 12 casos (2.1\%) fuera de [-0.288, 0.786] por regla IQR. & mayor en maligno; |d|=1.81; asimetría marcada. \\
worst concave points & numérica continua & Peor valor de cantidad relativa de puntos cóncavos del contorno. & 0 a 0.291 & índice adimensional & Muy alta & 0 casos (0.0\%) fuera de [-0.0798, 0.306] por regla IQR. & mayor en maligno; |d|=2.69; distribución relativamente estable. \\
worst symmetry & numérica continua & Peor valor de grado de simetría de la estructura celular. & 0.1565 a 0.6638 & índice adimensional & Media & 23 casos (4.0\%) fuera de [0.149, 0.419] por regla IQR. & mayor en maligno; |d|=0.95; asimetría marcada. \\
worst fractal dimension & numérica continua & Peor valor de aproximación de complejidad del contorno celular. & 0.05504 a 0.2075 & índice adimensional & Media & 24 casos (4.2\%) fuera de [0.0405, 0.123] por regla IQR. & mayor en maligno; |d|=0.71; asimetría marcada. \\
\end{longtable}
\normalsize

\end{landscape}


\section{Anexo B. Estadística descriptiva completa}
Las Tablas~\ref{tab:descriptive-position} y~\ref{tab:descriptive-signal} presentan los estadísticos de posición, dispersión, asimetría, atípicos IQR y señal univariada conservados desde el reporte HTML.

\begin{landscape}

\tiny
\setlength{\tabcolsep}{1.5pt}
\renewcommand{\arraystretch}{1.16}
\begin{longtable}{>{\raggedright\arraybackslash}p{0.148\linewidth} >{\raggedright\arraybackslash}p{0.056\linewidth} >{\raggedright\arraybackslash}p{0.074\linewidth} >{\raggedright\arraybackslash}p{0.083\linewidth} >{\raggedright\arraybackslash}p{0.074\linewidth} >{\raggedright\arraybackslash}p{0.074\linewidth} >{\raggedright\arraybackslash}p{0.074\linewidth} >{\raggedright\arraybackslash}p{0.074\linewidth} >{\raggedright\arraybackslash}p{0.074\linewidth} >{\raggedright\arraybackslash}p{0.074\linewidth} >{\raggedright\arraybackslash}p{0.074\linewidth}}
\caption{Estadísticos descriptivos de posición y dispersión}\label{tab:descriptive-position}\\
\toprule
Variable & n & Media & Desv. estándar & Mínimo & P1 & P25 & Mediana & P75 & P99 & Máximo \\
\midrule
\endfirsthead
\caption[]{Estadísticos descriptivos de posición y dispersión (continuación)}\\
\toprule
Variable & n & Media & Desv. estándar & Mínimo & P1 & P25 & Mediana & P75 & P99 & Máximo \\
\midrule
\endhead
\midrule
\multicolumn{11}{r}{Continúa en la siguiente página}\\
\midrule
\endfoot
\bottomrule
\endlastfoot
mean radius & 569.0 & 14.13 & 3.524 & 6.981 & 8.4584 & 11.70 & 13.37 & 15.78 & 24.37 & 28.11 \\
mean texture & 569.0 & 19.29 & 4.301 & 9.71 & 10.93 & 16.17 & 18.84 & 21.80 & 30.65 & 39.28 \\
mean perimeter & 569.0 & 91.97 & 24.30 & 43.79 & 53.83 & 75.17 & 86.24 & 104.1 & 165.7 & 188.5 \\
mean area & 569.0 & 654.9 & 351.9 & 143.5 & 215.7 & 420.3 & 551.1 & 782.7 & 1,786.6 & 2,501.0 \\
mean smoothness & 569.0 & 0.0964 & 0.0141 & 0.0526 & 0.0687 & 0.0864 & 0.0959 & 0.1053 & 0.1329 & 0.1634 \\
mean compactness & 569.0 & 0.1043 & 0.0528 & 0.0194 & 0.0334 & 0.0649 & 0.0926 & 0.1304 & 0.2772 & 0.3454 \\
mean concavity & 569.0 & 0.0888 & 0.0797 & 0 & 0 & 0.0296 & 0.0615 & 0.1307 & 0.3517 & 0.4268 \\
mean concave points & 569.0 & 0.0489 & 0.0388 & 0 & 0 & 0.0203 & 0.0335 & 0.074 & 0.1642 & 0.2012 \\
mean symmetry & 569.0 & 0.1812 & 0.0274 & 0.106 & 0.1295 & 0.1619 & 0.1792 & 0.1957 & 0.2596 & 0.304 \\
mean fractal dimension & 569.0 & 0.0628 & 0.0071 & 0.05 & 0.0515 & 0.0577 & 0.0615 & 0.0661 & 0.0854 & 0.0974 \\
radius error & 569.0 & 0.4052 & 0.2773 & 0.1115 & 0.1197 & 0.2324 & 0.3242 & 0.4789 & 1.2913 & 2.873 \\
texture error & 569.0 & 1.2169 & 0.5516 & 0.3602 & 0.4105 & 0.8339 & 1.108 & 1.474 & 2.9154 & 4.885 \\
perimeter error & 569.0 & 2.8661 & 2.0219 & 0.757 & 0.9532 & 1.606 & 2.287 & 3.357 & 9.69 & 21.98 \\
area error & 569.0 & 40.34 & 45.49 & 6.802 & 8.5144 & 17.85 & 24.53 & 45.19 & 177.7 & 542.2 \\
smoothness error & 569.0 & 0.007 & 0.003 & 0.0017 & 0.0031 & 0.0052 & 0.0064 & 0.0081 & 0.0173 & 0.0311 \\
compactness error & 569.0 & 0.0255 & 0.0179 & 0.0023 & 0.0047 & 0.0131 & 0.0204 & 0.0324 & 0.0899 & 0.1354 \\
concavity error & 569.0 & 0.0319 & 0.0302 & 0 & 0 & 0.0151 & 0.0259 & 0.042 & 0.1223 & 0.396 \\
concave points error & 569.0 & 0.0118 & 0.0062 & 0 & 0 & 0.0076 & 0.0109 & 0.0147 & 0.0312 & 0.0528 \\
symmetry error & 569.0 & 0.0205 & 0.0083 & 0.0079 & 0.0105 & 0.0152 & 0.0187 & 0.0235 & 0.0522 & 0.079 \\
fractal dimension error & 569.0 & 0.0038 & 0.0026 & 0.0009 & 0.0011 & 0.0022 & 0.0032 & 0.0046 & 0.0126 & 0.0298 \\
worst radius & 569.0 & 16.27 & 4.8332 & 7.93 & 9.2076 & 13.01 & 14.97 & 18.79 & 30.76 & 36.04 \\
worst texture & 569.0 & 25.68 & 6.1463 & 12.02 & 15.20 & 21.08 & 25.41 & 29.72 & 41.80 & 49.54 \\
worst perimeter & 569.0 & 107.3 & 33.60 & 50.41 & 58.27 & 84.11 & 97.66 & 125.4 & 208.3 & 251.2 \\
worst area & 569.0 & 880.6 & 569.4 & 185.2 & 256.2 & 515.3 & 686.5 & 1,084.0 & 2,918.2 & 4,254.0 \\
worst smoothness & 569.0 & 0.1324 & 0.0228 & 0.0712 & 0.0879 & 0.1166 & 0.1313 & 0.146 & 0.1889 & 0.2226 \\
worst compactness & 569.0 & 0.2543 & 0.1573 & 0.0273 & 0.0501 & 0.1472 & 0.2119 & 0.3391 & 0.7786 & 1.058 \\
worst concavity & 569.0 & 0.2722 & 0.2086 & 0 & 0 & 0.1145 & 0.2267 & 0.3829 & 0.9024 & 1.252 \\
worst concave points & 569.0 & 0.1146 & 0.0657 & 0 & 0 & 0.0649 & 0.0999 & 0.1614 & 0.2692 & 0.291 \\
worst symmetry & 569.0 & 0.2901 & 0.0619 & 0.1565 & 0.176 & 0.2504 & 0.2822 & 0.3179 & 0.4869 & 0.6638 \\
worst fractal dimension & 569.0 & 0.0839 & 0.0181 & 0.055 & 0.0586 & 0.0715 & 0.08 & 0.0921 & 0.1406 & 0.2075 \\
\end{longtable}
\normalsize

\end{landscape}


\begin{landscape}

\scriptsize
\setlength{\tabcolsep}{1.5pt}
\renewcommand{\arraystretch}{1.16}
\begin{longtable}{>{\raggedright\arraybackslash}p{0.235\linewidth} >{\raggedright\arraybackslash}p{0.117\linewidth} >{\raggedright\arraybackslash}p{0.098\linewidth} >{\raggedright\arraybackslash}p{0.176\linewidth} >{\raggedright\arraybackslash}p{0.117\linewidth} >{\raggedright\arraybackslash}p{0.137\linewidth}}
\caption{Asimetría, señal univariada y atípicos IQR}\label{tab:descriptive-signal}\\
\toprule
Variable & Asimetría & |d| & Dirección & Atípicos IQR & Atípicos IQR (\%) \\
\midrule
\endfirsthead
\caption[]{Asimetría, señal univariada y atípicos IQR (continuación)}\\
\toprule
Variable & Asimetría & |d| & Dirección & Atípicos IQR & Atípicos IQR (\%) \\
\midrule
\endhead
\midrule
\multicolumn{6}{r}{Continúa en la siguiente página}\\
\midrule
\endfoot
\bottomrule
\endlastfoot
mean radius & 0.9424 & 2.2055 & Mayor en maligno & 14.00 & 2.4605 \\
mean texture & 0.6504 & 0.9423 & Mayor en maligno & 7 & 1.2302 \\
mean perimeter & 0.9907 & 2.2895 & Mayor en maligno & 13.00 & 2.2847 \\
mean area & 1.6457 & 2.0757 & Mayor en maligno & 25.00 & 4.3937 \\
mean smoothness & 0.4563 & 0.793 & Mayor en maligno & 6 & 1.0545 \\
mean compactness & 1.1901 & 1.5346 & Mayor en maligno & 16.00 & 2.812 \\
mean concavity & 1.4012 & 2.0033 & Mayor en maligno & 18.00 & 3.1634 \\
mean concave points & 1.1712 & 2.5452 & Mayor en maligno & 10.00 & 1.7575 \\
mean symmetry & 0.7256 & 0.723 & Mayor en maligno & 15.00 & 2.6362 \\
mean fractal dimension & 1.3045 & 0.0265 & Mayor en benigno & 15.00 & 2.6362 \\
radius error & 3.0886 & 1.4217 & Mayor en maligno & 38.00 & 6.6784 \\
texture error & 1.6464 & 0.0171 & Mayor en benigno & 20.00 & 3.5149 \\
perimeter error & 3.4436 & 1.3816 & Mayor en maligno & 38.00 & 6.6784 \\
area error & 5.4472 & 1.3534 & Mayor en maligno & 65.00 & 11.42 \\
smoothness error & 2.3145 & 0.1387 & Mayor en benigno & 30.00 & 5.2724 \\
compactness error & 1.9022 & 0.6327 & Mayor en maligno & 28.00 & 4.9209 \\
concavity error & 5.1105 & 0.5416 & Mayor en maligno & 22.00 & 3.8664 \\
concave points error & 1.4447 & 0.9228 & Mayor en maligno & 19.00 & 3.3392 \\
symmetry error & 2.1951 & 0.0135 & Mayor en benigno & 27.00 & 4.7452 \\
fractal dimension error & 3.924 & 0.1615 & Mayor en maligno & 28.00 & 4.9209 \\
worst radius & 1.1031 & 2.5439 & Mayor en maligno & 17.00 & 2.9877 \\
worst texture & 0.4983 & 1.0605 & Mayor en maligno & 5 & 0.8787 \\
worst perimeter & 1.1282 & 2.5982 & Mayor en maligno & 15.00 & 2.6362 \\
worst area & 1.8594 & 2.2302 & Mayor en maligno & 35.00 & 6.1511 \\
worst smoothness & 0.4154 & 0.9596 & Mayor en maligno & 7 & 1.2302 \\
worst compactness & 1.4736 & 1.5126 & Mayor en maligno & 16.00 & 2.812 \\
worst concavity & 1.1502 & 1.8119 & Mayor en maligno & 12.00 & 2.109 \\
worst concave points & 0.4926 & 2.6926 & Mayor en maligno & 0 & 0 \\
worst symmetry & 1.4339 & 0.9453 & Mayor en maligno & 23.00 & 4.0422 \\
worst fractal dimension & 1.6626 & 0.7068 & Mayor en maligno & 24.00 & 4.2179 \\
\end{longtable}
\normalsize

\end{landscape}


\section{Anexo C. Reproducibilidad}
El proyecto se reproduce desde la raíz del repositorio con:

\begin{verbatim}
python src/build_project.py
\end{verbatim}

El PDF se compone a partir de los artefactos finales con:

\begin{verbatim}
python src/build_pdf_report.py
\end{verbatim}

Durante la generación de este documento no se recalcularon estadísticas ni se regeneraron gráficos analíticos. Solo se transformaron los artefactos existentes en una versión académica en PDF.

\end{document}
